\section{Η Εποχή του \ENG{AICC} (1988-2004)}
\label{sec:aicc_era}

Η \textit{\ENG{Aviation} \ENG{Industry} \ENG{Computer-Based} \ENG{Training} \ENG{Committee}} (\ENG{AICC}) αποτελεί την πρώτη οργανωμένη προσπάθεια προτυποποίησης στον τομέα του ηλεκτρονικού εκπαιδευτικού περιεχομένου. Ιδρύθηκε το 1988 ως διεθνής ένωση εταιρειών της αεροπορικής βιομηχανίας με στόχο την ανάπτυξη κατευθυντήριων γραμμών για την παραγωγή, την παράδοση και την αξιολόγηση εκπαιδευτικού υλικού βασισμένου σε υπολογιστές~\cite{wiki_aicc,growtheng2024aicc}. Η δημιουργία του \ENG{AICC} ήταν απάντηση στο υψηλό κόστος ανάπτυξης και συντήρησης εκπαιδευτικών προγραμμάτων στην αεροπορική βιομηχανία, όπου η τεχνική πολυπλοκότητα και οι αυστηρές κανονιστικές απαιτήσεις καθιστούσαν την εκπαίδευση κρίσιμο παράγοντα ασφάλειας και αποδοτικότητας.

\subsection{Ίδρυση και Στόχοι του \ENG{AICC}}
\label{subsec:aicc_foundation}

Τα μέλη του \ENG{AICC} περιλάμβαναν μεγάλες αεροπορικές εταιρείες όπως οι \ENG{Boeing}, \ENG{Airbus}, \ENG{American} \ENG{Airlines} και \ENG{Lufthansa}, καθώς και κατασκευαστές αεροσκαφών, προμηθευτές εκπαιδευτικής τεχνολογίας και ρυθμιστικές αρχές. Ο οργανισμός λειτουργούσε με βάση την αρχή της εθελοντικής συμμετοχής και της συναινετικής ανάπτυξης προδιαγραφών, χωρίς νομική εξαναγκαστική ισχύ αλλά με ισχυρή βιομηχανική αποδοχή~\cite{elearningindustry2024standards}.

Κεντρικός στόχος του \ENG{AICC} ήταν η μείωση του κόστους ανάπτυξης εκπαιδευτικού περιεχομένου μέσω της επαναχρησιμοποίησης και της διαλειτουργικότητας. Παράλληλα επιδίωξε τη δημιουργία κοινών τεχνικών κατευθυντήριων γραμμών για την υλοποίηση συστημάτων \ENG{CBT}, τη διασφάλιση συμβατότητας εκπαιδευτικού περιεχομένου μεταξύ διαφορετικών πλατφορμών και προμηθευτών, καθώς και την καθιέρωση μηχανισμών παρακολούθησης και αναφοράς της προόδου των εκπαιδευομένων.

Στο πλαίσιο αυτό, το \ENG{AICC} εξέδωσε μια σειρά τεχνικών συστάσεων και κατευθυντήριων γραμμών (\textit{\ENG{guidelines}}), γνωστών ως \textit{\ENG{AICC} \ENG{Guidelines} \ENG{and} \ENG{Recommendations}} (\textit{\ENG{AGRs}}). Οι προδιαγραφές αυτές κάλυπταν σημαντικές πτυχές των συστημάτων \ENG{E-Learning}, όπως τη δομή του περιεχομένου, τους μηχανισμούς επικοινωνίας μεταξύ εκπαιδευτικού υλικού και \ENG{LMS}, καθώς και τη χρήση μεταδεδομένων για την περιγραφή εκπαιδευτικών πόρων~\cite{paradiso2024aicc,teachfloor2024aicc}.

\subsection{\ENG{AICC} \ENG{CMI} \ENG{Guidelines} (\ENG{CMI001})}
\label{subsec:aicc_cmi}

Ένα από τα σημαντικότερα επιτεύγματα του \ENG{AICC} ήταν η ανάπτυξη των προδιαγραφών \textit{\ENG{Computer} \ENG{Managed} \ENG{Instruction}} (\ENG{CMI}), γνωστών ως \ENG{CMI001}. Το σύνολο αυτό προδιαγραφών όρισε ένα τυποποιημένο μοντέλο δεδομένων (\textit{\ENG{data} \ENG{model}}) και ένα σύνολο πρωτοκόλλων επικοινωνίας για την ανταλλαγή πληροφοριών μεταξύ εκπαιδευτικού περιεχομένου (\textit{\ENG{courseware}}) και συστημάτων διαχείρισης μάθησης~\cite{tlaworks2024evolution}.

\subsubsection{\ENG{HACP}: \ENG{HTTP-based} \ENG{AICC} \ENG{Communication} \ENG{Protocol}}

Στον πυρήνα της τεχνικής υλοποίησης βρισκόταν το \textit{\ENG{AICC} \ENG{HTTP-based} \ENG{Communication} \ENG{Protocol}} (\ENG{HACP}). Το πρωτόκολλο αυτό επέτρεπε στο εκπαιδευτικό περιεχόμενο να επικοινωνεί με το \ENG{LMS} μέσω αιτημάτων \ENG{HTTP}, ανεξάρτητα από την πλατφόρμα ή τον φυλλομετρητή (\textit{\ENG{browser}}) που χρησιμοποιούνταν. Για την εποχή του, η προσέγγιση αυτή ήταν ιδιαίτερα καινοτόμος, καθώς επέτρεπε την παροχή εκπαιδευτικού περιεχομένου μέσω \ENG{web} \ENG{browsers} και την κεντρική παρακολούθηση της προόδου των εκπαιδευομένων μέσω δικτύου.

Το \ENG{HACP} βασιζόταν σε έναν απλό μηχανισμό ανταλλαγής δεδομένων μέσω \textit{\ENG{name-value} \ENG{pairs}}. Το εκπαιδευτικό περιεχόμενο έστελνε αιτήματα \ENG{HTTP} \ENG{GET} ή \ENG{POST} προς το \ENG{LMS}, περιλαμβάνοντας παραμέτρους όπως το \texttt{\ENG{aicc}\_\ENG{sid}} για το \ENG{session} \ENG{ID}, το \texttt{\ENG{aicc}\_\ENG{url}} για τη διεύθυνση του \ENG{LMS} και το \texttt{\ENG{command}} για τον καθορισμό της ενέργειας που ζητείται. Το \ENG{LMS} απαντούσε με δομημένες απαντήσεις που περιείχαν πληροφορίες προόδου ή επιβεβαιώσεις εκτέλεσης των εντολών.

\subsubsection{\ENG{Data} \ENG{Model} \ENG{Elements}}

Το μοντέλο δεδομένων του \ENG{AICC} όριζε ένα σύνολο στοιχείων (\textit{\ENG{data} \ENG{elements}}) που μπορούσαν να ανταλλάσσονται μεταξύ του εκπαιδευτικού περιεχομένου και του \ENG{LMS}. Τα στοιχεία αυτά περιλάμβαναν πληροφορίες σχετικές με την ταυτότητα του χρήστη, την κατάσταση ολοκλήρωσης του μαθήματος, τα αποτελέσματα αξιολόγησης (\ENG{Quiz}) και τη θέση προόδου μέσα στο εκπαιδευτικό υλικό.

\begin{table}[H]
\centering
\caption{Στοιχεία του \ENG{AICC} \ENG{CMI} \ENG{data} \ENG{model}}
\label{tab:aicc_cmi_data_elements}
\small
\setlength{\tabcolsep}{5pt}
\renewcommand{\arraystretch}{1.25}
\begin{chaptertablebox}
\begin{tabular}{|l|p{8cm}|}
\hline
\textbf{Στοιχείο} & \textbf{Περιγραφή} \\
\hline
\hline
\ENG{Student} \ENG{Identification} & Στοιχεία ταυτοποίησης του εκπαιδευόμενου (π.χ. \texttt{\ENG{student}\_\ENG{id}}, \texttt{\ENG{student}\_\ENG{name}}). \\
\ENG{Lesson} \ENG{Status} & Κατάσταση ολοκλήρωσης του μαθήματος (π.χ. \texttt{\ENG{incomplete}}, \texttt{\ENG{completed}}, \texttt{\ENG{passed}}, \texttt{\ENG{failed}}). \\
\ENG{Score} & Βαθμολογία του εκπαιδευόμενου (π.χ. \texttt{\ENG{score}} σε ποσοστό ή απόλυτη τιμή). \\
\ENG{Time} & Χρονικές πληροφορίες (π.χ. \texttt{\ENG{time}}, συνολικός χρόνος μελέτης). \\
\ENG{Lesson} \ENG{Location} & Η τελευταία θέση στην οποία βρισκόταν ο εκπαιδευόμενος (\ENG{bookmarking}). \\
\ENG{Suspend} \ENG{Data} & Δεδομένα που επιτρέπουν την αναστολή και συνέχιση της εκπαιδευτικής συνεδρίας. \\
\hline
\end{tabular}
\end{chaptertablebox}
\end{table}


Το \ENG{data} \ENG{model} του \ENG{AICC} αποτέλεσε τη βάση για πολλά μεταγενέστερα πρότυπα, ιδίως για το \ENG{SCORM} 1.2, το οποίο υιοθέτησε αρκετές από τις βασικές έννοιες και τη δομή του \ENG{AICC}~\cite{elearningchef2024comparison}.

\subsection{Αρχεία Περιγραφής Μαθημάτων (\ENG{Course} \ENG{Structure} \ENG{Files})}

Το \ENG{AICC} καθόρισε επίσης προδιαγραφές για αρχεία που περιέγραφαν τη δομή των εκπαιδευτικών μαθημάτων. Συγκεκριμένα, ορίστηκαν διαφορετικοί τύποι αρχείων για την περιγραφή της συνολικής δομής ενός μαθήματος, των επιμέρους εκπαιδευτικών μονάδων, των μεταδεδομένων και της ιεραρχίας των δραστηριοτήτων. Τα αρχεία αυτά περιλάμβαναν το \ENG{Course} \ENG{Description} \ENG{File} (.\ENG{crs}), το οποίο περιέγραφε τη συνολική δομή του μαθήματος, το \ENG{Assignable} \ENG{Unit} \ENG{File} (.\ENG{au}), το οποίο όριζε τις επιμέρους μονάδες που μπορούσαν να ανατεθούν σε εκπαιδευόμενους, το \ENG{Descriptor} \ENG{File} (.\ENG{des}) που περιείχε μεταδεδομένα, καθώς και το \ENG{Course} \ENG{Structure} \ENG{File} (.\ENG{cst}), το οποίο καθόριζε την ιεραρχία και τη σειρά των εκπαιδευτικών μονάδων.

Τα αρχεία αυτά ήταν σε μορφή απλού κειμένου (\ENG{plain} \ENG{text}) και χρησιμοποιούσαν μια απλή σύνταξη \textit{\ENG{name}=\ENG{value}}, γεγονός που τα καθιστούσε εύκολα αναγνώσιμα και επεξεργάσιμα.

\subsection{Κληρονομιά και Επιρροή στο \ENG{SCORM}}
\label{subsec:aicc_legacy}

Η επιρροή του \ENG{AICC} στην εξέλιξη των \ENG{E-Learning} προτύπων υπήρξε ιδιαίτερα σημαντική. Πολλές από τις βασικές έννοιες που εισήγαγε υιοθετήθηκαν αργότερα από την πρωτοβουλία \textit{\ENG{Advanced} \ENG{Distributed} \ENG{Learning}} (\ENG{ADL}) κατά την ανάπτυξη του \textit{\ENG{Sharable} \ENG{Content} \ENG{Object} \ENG{Reference} \ENG{Model}} (\ENG{SCORM}). Το μοντέλο δεδομένων του \ENG{SCORM} 1.2 βασίστηκε σε μεγάλο βαθμό στο \ENG{AICC} \ENG{CMI} \ENG{data} \ENG{model}, εμπλουτισμένο με νέα στοιχεία και βελτιώσεις στη δομή. Παράλληλα, έννοιες όπως το \ENG{bookmarking}, δηλαδή η αποθήκευση της τρέχουσας θέσης του εκπαιδευομένου μέσα στο μάθημα, προέρχονται άμεσα από τις προσεγγίσεις του \ENG{AICC}. Επιπλέον, το πρωτόκολλο \ENG{HACP} επηρέασε τη σχεδίαση του \ENG{SCORM} \ENG{Run-Time} \ENG{Environment} (\ENG{RTE}), αν και το \ENG{SCORM} υιοθέτησε διαφορετική τεχνική προσέγγιση βασισμένη σε \ENG{JavaScript} \ENG{API} αντί για \ENG{HTTP} αιτήματα.

Παρά την αρχική επιτυχία του, το \ENG{AICC} παρουσίασε ορισμένους περιορισμούς. Η υποστήριξη για σύνθετες δομές μαθημάτων και προηγμένους μηχανισμούς αλληλουχίας (\textit{\ENG{sequencing}}) ήταν περιορισμένη, ενώ δεν υπήρχε ένας πλήρως τυποποιημένος μηχανισμός συσκευασίας και διανομής του περιεχομένου. Επιπλέον, η επικοινωνία μέσω \ENG{HACP} απαιτούσε σύνδεση με διακομιστή μέσω \ENG{HTTP}, γεγονός που περιόριζε τη χρήση σε περιβάλλοντα χωρίς δικτυακή σύνδεση. Η υποστήριξη μεταδεδομένων παρέμεινε επίσης περιορισμένη, καθώς το πρότυπο δεν ενσωμάτωσε πλήρως νεότερα σχήματα μεταδεδομένων, όπως το \ENG{IEEE} \ENG{Learning} \ENG{Object} \ENG{Metadata} (\ENG{LOM}).

Παρά τους περιορισμούς αυτούς, το \ENG{AICC} παρέμεινε ευρέως χρησιμοποιούμενο μέχρι τις αρχές της δεκαετίας του 2000, και πολλά \ENG{LMS} συνέχισαν να υποστηρίζουν \ENG{AICC} περιεχόμενο ακόμη και μετά την εμφάνιση του \ENG{SCORM}. Η έμφαση που έδωσε στη διαλειτουργικότητα και η πρακτική προσέγγιση που υιοθέτησε στη διαδικασία προτυποποίησης συνέβαλαν καθοριστικά στη διαμόρφωση του οικοσυστήματος των προτύπων \ENG{E-Learning}.

Το 2014 ο οργανισμός \ENG{AICC} ανακοίνωσε την επίσημη διάλυσή του, αναγνωρίζοντας ότι νεότεροι οργανισμοί προτυποποίησης όπως το \ENG{ADL} και το \ENG{IMS} \ENG{Global} είχαν πλέον αναλάβει τον ηγετικό ρόλο στην εξέλιξη των προτύπων \ENG{E-Learning}~\cite{opigno2024standards}. Παρ' όλα αυτά, η ιστορική του συμβολή παραμένει σημαντική, καθώς πολλές από τις βασικές έννοιες που εισήγαγε εξακολουθούν να επηρεάζουν τα σύγχρονα πρότυπα.

\begin{figure}[H]
  \centering
  % FIGURE 2.3: AICC HACP Protocol (HTTP-based AICC Communication)
% Sequence diagram showing LMS-Content communication
\begin{chapterfigurebox}
\begin{tikzpicture}[
  scale=0.9,
  every node/.style={font=\small},
  entity/.style={rectangle, draw, thick, minimum width=3cm, minimum height=1.2cm, align=center, fill=white},
  message/.style={-latex, thick},
  return/.style={-latex, thick, dashed}
]

% Entities
\node[entity, fill=aiccblue!20] at (0, 8) (lms) {\textbf{\ENG{LMS}}\\(\ENG{Learning} \ENG{Management} \ENG{System})};
\node[entity, fill=aiccblue!40] at (8, 8) (content) {\textbf{\ENG{AICC} \ENG{Content}}\\(\ENG{AU} - \ENG{Assignable} \ENG{Unit})};

% Lifelines
\draw[thick, dashed] (lms.south) -- (0, -1);
\draw[thick, dashed] (content.south) -- (8, -1);

% Step 1: Launch
\draw[message, aiccblue] (0, 6.5) -- (8, 6.5) node[midway, above, font=\footnotesize] {1. \ENG{Launch} \ENG{AU} \ENG{with} \ENG{parameters}};
\node[right, font=\tiny, align=left, text width=4cm] at (8.5, 6.5) {
  \texttt{\ENG{AICC}\_\ENG{URL}=\ENG{http}://...}\\
  \texttt{\ENG{AICC}\_\ENG{SID}=\ENG{session123}}
};

% Step 2: GetParam request
\draw[message, xapiorange] (8, 5.5) -- (0, 5.5) node[midway, above, font=\footnotesize] {2. \ENG{HACP}\_\ENG{GetParam}()};
\node[left, font=\tiny, align=left] at (-0.5, 5.5) {
  \ENG{HTTP} \ENG{GET}:\\
  \texttt{?\ENG{command}=\ENG{GetParam}}
};

% Step 3: GetParam response
\draw[return, xapiorange] (0, 4.5) -- (8, 4.5) node[midway, above, font=\footnotesize] {3. \ENG{Return} \ENG{learner} \ENG{data}};
\node[right, font=\tiny, align=left, text width=4.5cm] at (8.5, 4.5) {
  \texttt{[\ENG{Core}]}\\
  \texttt{\ENG{Student}\_\ENG{ID}=\ENG{john123}}\\
  \texttt{\ENG{Student}\_\ENG{Name}=\ENG{John} \ENG{Doe}}\\
  \texttt{\ENG{Lesson}\_\ENG{Status}=\ENG{incomplete}}
};

% Step 4: Learner interaction (internal)
\draw[->, thick, gray] (8, 3.8) -- (8, 3.2);
\node[right, font=\footnotesize, gray] at (8.5, 3.5) {\ENG{Learner} \ENG{interacts} \ENG{with} \ENG{content}};

% Step 5: PutParam request
\draw[message, scormgreen] (8, 2.5) -- (0, 2.5) node[midway, above, font=\footnotesize] {4. \ENG{HACP}\_\ENG{PutParam}()};
\node[left, font=\tiny, align=left, text width=4cm] at (-0.5, 2.5) {
  \ENG{HTTP} \ENG{POST}:\\
  \texttt{?\ENG{command}=\ENG{PutParam}}\\
  \texttt{\&\ENG{lesson}\_\ENG{status}=\ENG{completed}}\\
  \texttt{\&\ENG{score}=85}
};

% Step 6: PutParam acknowledgment
\draw[return, scormgreen] (0, 1.5) -- (8, 1.5) node[midway, above, font=\footnotesize] {5. \ENG{ACK} (\ENG{error}=0)};

% Step 7: More interactions (optional)
\draw[->, thick, gray, dashed] (8, 1.0) -- (8, 0.5);
\node[right, font=\footnotesize, gray] at (8.5, 0.75) {\ENG{Additional} \ENG{PutParam} \ENG{calls...}};

% Step 8: ExitAU
\draw[message, problemred] (8, 0) -- (0, 0) node[midway, above, font=\footnotesize] {6. \ENG{HACP}\_\ENG{ExitAU}()};
\node[left, font=\tiny, align=left] at (-0.5, 0) {
  \ENG{HTTP} \ENG{GET}:\\
  \texttt{?\ENG{command}=\ENG{ExitAU}}
};

% Legend
\node[below, font=\tiny, align=left] at (4, -1.5) {
  \textbf{\ENG{HACP} \ENG{Commands}:}\\
  \textcolor{xapiorange}{\textbullet} \textbf{\ENG{GetParam}:} \ENG{Retrieve} \ENG{initial} \ENG{learner/course} \ENG{data} \ENG{from} \ENG{LMS}\\
  \textcolor{scormgreen}{\textbullet} \textbf{\ENG{PutParam}:} \ENG{Send} \ENG{tracking} \ENG{data} (\ENG{score}, \ENG{status}, \ENG{time}) \ENG{to} \ENG{LMS}\\
  \textcolor{problemred}{\textbullet} \textbf{\ENG{ExitAU}:} \ENG{Notify} \ENG{LMS} \ENG{of} \ENG{content} \ENG{completion/exit}
};

% AICC data model note
\node[below, font=\scriptsize, align=center, text width=10cm] at (4, -3.2) {
  \textbf{\ENG{AICC} \ENG{Data} \ENG{Model}:} \ENG{Core}, \ENG{Core}\_\ENG{Lesson}, \ENG{Core}\_\ENG{Student}, \ENG{Core}\_\ENG{Vendor}, \ENG{Evaluation}, \ENG{Student}\_\ENG{Demographics}\\
  \textbf{\ENG{Transport}:} \ENG{HTTP} (\ENG{HACP} = \ENG{HTTP-based} \ENG{AICC} \ENG{Communication} \ENG{Protocol})\\
  \textbf{\ENG{Format}:} \ENG{Key-value} \ENG{pairs} (\ENG{INI-style}), \textbf{\ENG{NOT}} \ENG{XML}
};

% Advantages box
\node[rectangle, draw, thick, fill=solutiongreen!10, align=left, font=\scriptsize, text width=5cm] at (-3, 3.5) {
  \textbf{\ENG{AICC} \ENG{Advantages} (1988):}\\
  \textbullet~\ENG{First} \ENG{industry} \ENG{standard}\\
  \textbullet~\ENG{HTTP-based} (\ENG{portable})\\
  \textbullet~\ENG{Offline} \ENG{support} (\ENG{AICC}\_\ENG{URL})\\
  \textbullet~\ENG{Simple} \ENG{key-value} \ENG{format}
};

\end{tikzpicture}
\end{chapterfigurebox}

  \caption{\ENG{AICC} \ENG{HACP} \ENG{Protocol} - \ENG{HTTP-based} \ENG{Communication} \ENG{Flow}}
  \label{fig:aicc_hacp}
\end{figure}

\setlength{\LTleft}{0pt plus 1fil}
\setlength{\LTright}{0pt plus 1fil}
\small
\setlength{\tabcolsep}{5pt}
\renewcommand{\arraystretch}{1.25}
\begin{longtable}{@{}p{3.5cm}p{5cm}p{5.5cm}@{}}
\caption{Σύγκριση \ENG{AICC} και \ENG{SCORM} - Χαρακτηριστικά και Κληρονομιά}
\label{tab:aicc_vs_scorm}\\
\toprule
\textbf{Χαρακτηριστικό} & \textbf{\ENG{AICC} (1988-2014)} & \textbf{\ENG{SCORM} 1.2/2004 (2001-σήμερα)} \\
\midrule
\endfirsthead

\caption[]{Σύγκριση \ENG{AICC} και \ENG{SCORM} - Χαρακτηριστικά και Κληρονομιά (συνέχεια)}\\
\toprule
\textbf{Χαρακτηριστικό} & \textbf{\ENG{AICC} (1988-2014)} & \textbf{\ENG{SCORM} 1.2/2004 (2001-σήμερα)} \\
\midrule
\endhead

\midrule
\multicolumn{3}{r}{\footnotesize Συνεχίζεται στην επόμενη σελίδα}\\
\endfoot

\bottomrule
\endlastfoot
\rowcolor{gray!10}
\textbf{Πρωτόκολλο Επικοινωνίας} & 
\ENG{HTTP-based} \ENG{AICC} \ENG{Communication} \ENG{Protocol} (\ENG{HACP})
\newline \textbullet{} \ENG{HTTP} \ENG{GET/POST} \ENG{requests}
\newline \textbullet{} \ENG{Name-value} \ENG{pairs}
\newline \textbullet{} \ENG{Server-dependent} &
\ENG{JavaScript} \ENG{API}
\newline \textbullet{} \ENG{Browser-based} (\ENG{window.API} \ENG{or} \ENG{API}\_1484\_11)
\newline \textbullet{} \ENG{Function} \ENG{calls} (\ENG{LMSInitialize}, \ENG{LMSGetValue}, \ENG{etc.})
\newline \textbullet{} \ENG{Client-side} \ENG{execution} \\
\midrule
\textbf{\ENG{Data} \ENG{Model}} & 
\ENG{CMI001} \ENG{specification}
\newline \textbullet{} \ENG{student}\_\ENG{id}, \ENG{student}\_\ENG{name}
\newline \textbullet{} \ENG{lesson}\_\ENG{status}
\newline \textbullet{} \ENG{score}, \ENG{time}
\newline \textbullet{} \ENG{lesson}\_\ENG{location}
\newline \textbullet{} \ENG{suspend}\_\ENG{data} &
\ENG{SCORM} 1.2: \ENG{CMI} \ENG{data} \ENG{model} (\ENG{extended} \ENG{from} \ENG{AICC})
\newline \ENG{SCORM} 2004: \ENG{CMI-5} \ENG{extended} \ENG{model}
\newline \textbullet{} \ENG{cmi.core.student}\_\ENG{id}
\newline \textbullet{} \ENG{cmi.core.lesson}\_\ENG{status}
\newline \textbullet{} \ENG{cmi.core.score.raw}
\newline \textbullet{} \ENG{cmi.suspend}\_\ENG{data} (4096 \ENG{chars} \ENG{in} 2004) \\
\midrule
\rowcolor{gray!10}
\textbf{\ENG{Content} \ENG{Packaging}} & 
\ENG{Course} \ENG{Structure} \ENG{Files}:
\newline \textbullet{} .\ENG{crs} (\ENG{Course} \ENG{Description})
\newline \textbullet{} .\ENG{au} (\ENG{Assignable} \ENG{Units})
\newline \textbullet{} .\ENG{des} (\ENG{Descriptor})
\newline \textbullet{} .\ENG{cst} (\ENG{Course} \ENG{Structure})
\newline \ENG{Plain} \ENG{text}, \ENG{name}=\ENG{value} \ENG{syntax} &
\ENG{IMS} \ENG{Content} \ENG{Packaging} (\ENG{manifest})
\newline \textbullet{} \ENG{imsmanifest.xml}
\newline \textbullet{} \ENG{Organizations}, \ENG{Resources}
\newline \textbullet{} \ENG{Metadata} (\ENG{IEEE} \ENG{LOM} \ENG{subset})
\newline \textbullet{} \ENG{ZIP} \ENG{package} (.\ENG{zip} \ENG{or} .\ENG{pif}) \\
\midrule
\textbf{\ENG{Sequencing} \& \ENG{Navigation}} & 
\textcolor{red}{Περιορισμένη υποστήριξη}
\newline \textbullet{} Βασική \ENG{linear} \ENG{navigation}
\newline \textbullet{} \ENG{Lesson}\_\ENG{location} για \ENG{bookmarking}
\newline \textbullet{} Δεν υπάρχει \ENG{sequencing} \ENG{logic} &
\ENG{SCORM} 1.2: Περιορισμένο (όπως \ENG{AICC})
\newline \textbf{\ENG{SCORM} 2004:} \ENG{IMS} \ENG{Simple} \ENG{Sequencing} (\ENG{SS})
\newline \textbullet{} \ENG{Sequencing} \ENG{rules}
\newline \textbullet{} \ENG{Objectives}, \ENG{rollup} \ENG{rules}
\newline \textbullet{} \ENG{Advanced} \ENG{navigation} \ENG{controls} \\
\midrule
\rowcolor{gray!10}
\textbf{\ENG{Metadata} \ENG{Support}} & 
\textcolor{orange}{\ENG{Minimal}}
\newline \textbullet{} \ENG{Basic} \ENG{descriptors} \ENG{in} .\ENG{des} \ENG{files}
\newline \textbullet{} Δεν ενσωματώνει \ENG{Dublin} \ENG{Core} ή \ENG{IEEE} \ENG{LOM} πλήρως &
\textcolor{ForestGreen}{Πλήρης}
\newline \textbullet{} \ENG{IEEE} \ENG{LOM} \ENG{subset} (\ENG{SCORM} 1.2)
\newline \textbullet{} \ENG{Full} \ENG{IEEE} \ENG{LOM} \ENG{support} (\ENG{SCORM} 2004)
\newline \textbullet{} \ENG{XML-based} \ENG{metadata} \ENG{in} \ENG{manifest} \\
\midrule
\textbf{\ENG{Offline} \ENG{Capability}} & 
\textcolor{red}{Περιορισμένη}
\newline Απαιτεί σύνδεση με \ENG{LMS} \ENG{server} για \ENG{HACP} \ENG{communication} &
\textcolor{ForestGreen}{Καλύτερη}
\newline Περιεχόμενο μπορεί να εκτελεστεί τοπικά
\newline \ENG{API} \ENG{calls} \ENG{buffered} \ENG{locally} (με περιορισμούς) \\
\midrule
\rowcolor{gray!10}
\textbf{\ENG{Adoption} \& \ENG{Legacy}} & 
\textbullet{} Πρώιμη υιοθέτηση (1990\ENG{s-early} 2000\ENG{s})
\newline \textbullet{} Κυρίως \ENG{aviation/corporate} \ENG{training}
\newline \textbullet{} Διάλυση οργανισμού 2014
\newline \textcolor{blue}{\textbf{\ENG{Legacy}:}} Βάση για \ENG{SCORM} \ENG{API} &
\textbullet{} Ευρεία υιοθέτηση (2001-σήμερα)
\newline \textbullet{} \ENG{SCORM} 1.2: \textasciitilde 95\% \ENG{LMS} \ENG{support}
\newline \textbullet{} \ENG{SCORM} 2004: Πολυπλοκότερο, αργότερη υιοθέτηση
\newline \textcolor{ForestGreen}{\textbf{\ENG{Current}:}} Παραμένει \ENG{active} \ENG{standard} \\
\midrule
\textbf{Τεχνικοί Περιορισμοί} & 
\textbullet{} Εξάρτηση από \ENG{HTTP} \ENG{server}
\newline \textbullet{} 128-\ENG{char} \ENG{data} \ENG{limits} (\ENG{suspend}\_\ENG{data})
\newline \textbullet{} Δεν υποστηρίζει προηγμένη \ENG{navigation}
\newline \textbullet{} Περιορισμένα \ENG{metadata} &
\ENG{SCORM} 1.2:
\newline \textbullet{} 4096-\ENG{char} \ENG{suspend}\_\ENG{data} \ENG{limit}
\newline \textbullet{} Δεν υποστηρίζει \ENG{sequencing}
\newline \ENG{SCORM} 2004:
\newline \textbullet{} Πολυπλοκότητα \ENG{implementation}
\newline \textbullet{} \ENG{LMS-centric} (δεν καλύπτει \ENG{mobile/social}) \\
\midrule
\rowcolor{gray!10}
\textbf{Κληρονομιά στα Σύγχρονα \ENG{Standards}} & 
\textbullet{} \ENG{HACP} \ENG{protocol} $\rightarrow$ \ENG{SCORM} \ENG{API} \ENG{design}
\newline \textbullet{} \ENG{CMI} \ENG{data} \ENG{model} $\rightarrow$ \ENG{SCORM} 1.2 \ENG{data} \ENG{model}
\newline \textbullet{} \ENG{Bookmarking} \ENG{concept} $\rightarrow$ \ENG{lesson}\_\ENG{location}
\newline \textbullet{} \ENG{Interoperability} \ENG{mindset} $\rightarrow$ όλα τα επόμενα \ENG{standards} &
\textbullet{} \ENG{SCORM} \ENG{API} $\rightarrow$ \ENG{cmi5} \ENG{API} (\ENG{xAPI} \ENG{wrapper})
\newline \textbullet{} \ENG{Content} \ENG{packaging} $\rightarrow$ \ENG{IMS} \ENG{Common} \ENG{Cartridge}
\newline \textbullet{} \ENG{Metadata} $\rightarrow$ \ENG{LRMI}, \ENG{schema.org} \ENG{extensions}
\newline \textbullet{} \ENG{Sequencing} \ENG{rules} $\rightarrow$ \ENG{xAPI} \ENG{profiles} (\ENG{adaptive} \ENG{learning}) \\
\end{longtable}
\normalsize









